\SourceLeaf{287}%
Werth $360^0$ desto näher kommt, je kleiner $\omega, \omega'$ sind; d.\,i. die Tangente an der Curve hat ihre Neigung vom Anfang bis zum Ende um $360^0$ geändert.

Schneidet die Curve sich einmal selbst, so bekommen diese Schlüsse eine Abänderung. Sind nemlich die Endpunkte von $s = \sigma$, $s' = \sigma'$ identisch, so leidet $(s', s)$ für diese Werthe eine Sprungsänderung: nemlich für $s = \sigma$, als constant betrachtet, ändert sich $(s', s)$, indem $s'$ den Werth $\sigma'$ erreicht, auf einmal um $\pm 180^0$; um zu entscheiden, welcher Werth hier gilt, müssen die Werthe für $s = \sigma \mp \omega$ betrachtet werden, wo $\omega$ unendlich klein ist; man sieht, dass, wie $s' - \sigma'$ von einem [negativen] Werthe [an] die immer kleiner werdenden negativen Werthe erhält und durch $0$ wieder wächst, $(s', \sigma - \omega)$ um eine Grösse zunimmt oder abnimmt, die unendlich wenig von $180^0$ verschieden ist, je nachdem man bei $s'$ von der innern auf die äussere Seite der Linie $s'$ übergeht oder umgekehrt; für $(s', \sigma + \omega)$ ist es hingegen umgekehrt; man sieht also, dass überhaupt für ein constantes $s'$, grösser als $\sigma'$, $(s', s)$ beim Durchgang durch den Werth $s = \sigma$ auf einmal um $360^0$ ab(zu)nimmt und dass die Änderung des $(s', \sigma)$, indem $s'$ durch den Werth $\sigma'$ geht, um $\pm 180^0$ eigentlich mit beiden Zeichen zugleich gilt, je nachdem man $s = \sigma$ als den letzten der \textls{bis} $\sigma$ oder als den ersten der \textls{von} $\sigma$ an weiter zu nehmenden Werthe ansieht. Das Solidum hat hier die Gestalt einer Wendeltreppe. Die Änderung der Richtung der Tangente wird nunmehro entweder $720^0$ oder $0^0$, man kann aber dieser Betrachtung ganz enthoben sein, wenn man die gegenwärtige Methode auf den Fall einschränkt, wo die Curve sich \textls{nicht} selbst schneidet. 

\begin{center}
[8.]
\end{center}

Überhaupt lässt sich in dieser Theorie vieles sehr vereinfachen, wenn man die Zerlegbarkeit der Polygone in Dreiecke voraussetzt. Diese aber lässt sich auf folgende Art beweisen.

Zuerst ist klar, dass jede Figur durch Eine Theilungslinie (gerade; kürzeste) sich in zwei zerlegen lässt, die zusammen einen Einwärts gehenden Winkel weniger haben. Durch wiederholte Theilung dieser Art wird also das Polygon in andere zerlegt, die gar keine einwärts gehende Winkel haben.

Es sei nun $P$ ein Polygon mit lauter auswärts gehenden Winkeln an den Punkten $(1), (2), (3) .\,.\,. (n)$; ich behaupte, die Linie $(1)(3)$ wird ganz im Innern des Polygons liegen.
